\subsection{Weak topologies on complex vector spaces}

The definition and results of II, § 6, Nos. 1 and 2 apply without change to complex vector spaces. If F and G are two complex vector spaces in duality by a bilinear form B, then the underlying spaces \(F_{0}\) and \(G_{0}\) are in duality by RB, and it follows from II, p.~61, formula (1) that the weak topologies \(\sigma(\mathrm{F},\mathrm{G})\) and \(\sigma(\mathrm{F}_{0},\mathrm{G}_{0})\) are identical.

DEFINITION 1. --- Let F and G be two complex vector spaces in duality. For any subset M of F, the polar of M in G, denoted by \(M^{\circ}\) , is the set of \(y \in G\) such that \(\mathcal{R}(\langle x, y \rangle) \geqslant -1\) for all \(x \in M\) .

If \(\mathbf{M}^{\circ}\) is the polar of \(\mathbf{M} \subset \mathbf{F}\) in \(\mathbf{G}\) then \((\lambda \mathbf{M})^{\circ} = \lambda^{-1}\mathbf{M}^{\circ}\) for all \(\lambda \in \mathbf{C}^{*}\) .

If M is a (complex) vector subspace of E, then \(M^{\circ}\) is a closed vector subspace (for \(\sigma(\mathbf{G},\mathbf{F})\) ), since the relation \(\mathcal{R}(\lambda\langle x,y\rangle)\geqslant-1\) for every scalar \(\lambda\in C\) implies \(\langle x,y\rangle=0\) ; again we say that \(M^{\circ}\) is the subspace of G orthogonal to M.

If M is a balanced set in F, then \(M^{\circ}\) is a balanced set in G; in this case \(M^{\circ}\) is the set of \(y \in G\) such that \(\left|\langle x, y \rangle\right| \leqslant 1\) for all \(x \in M\) ; for this relation is equivalent to \(\mathcal{R}(\langle \zeta x, y \rangle) \leqslant 1\) for all \(x \in M\) and all \(\zeta \in C\) such that \(|\zeta| = 1\) .

The results of II, p.~41 to 51 are also valid without restriction for complex vector spaces.

